\documentclass[11pt, a4paper]{article}

\usepackage[utf8]{inputenc}
\usepackage{amsmath, amssymb, amsthm}
\usepackage{geometry}
\geometry{margin=1in}
\usepackage{hyperref}
\usepackage{listings}
\usepackage{xcolor}

\lstdefinelanguage{GAP}{
  morekeywords={local, function, return, if, then, else, fi, for, in, do, od, while, repeat, until, rec, mod, and, or, not},
  sensitive=true,
  morecomment=[l]{\#},
  morestring=[b]",
}

\lstset{
  language=GAP,
  basicstyle=\ttfamily\small,
  keywordstyle=\color{blue}\bfseries,
  commentstyle=\color{green!60!black},
  stringstyle=\color{red},
  showstringspaces=false,
  frame=single,
  numbers=left,
  numberstyle=\tiny\color{gray},
  breaklines=true
}

\newtheorem{theorem}{Theorem}

\title{Algebraic Compression of Approximate Quantum Fourier Transform Amplitudes via Finite-Ring Dictionary Encoding}
\author{Sandeep Koranne}
\date{\today}

\begin{document}

\maketitle

\begin{abstract}
We introduce a finite-ring dictionary encoding for the exact analytical evaluation of discretized Coppersmith-type approximate quantum Fourier transform amplitudes. The method maps the phase function into the cyclotomic alphabet generated by a $2^K$-th root of unity, $\zeta_{2^K} = e^{2\pi i/2^K}$, with phase indices reduced modulo $2^K$ and represented through a compact $\mathbb{Z}_{2^K}$-valued dictionary. For binary superposition variables $x_j$ and fixed binary input $y$, the resulting phase is linear in the active variables, $Q(x,y) = \sum_j c_j(y)x_j$, where the coefficients $c_j(y)$ are obtained directly from the finite-range coupling structure. This permits the exponentially large amplitude sum $\sum_{x\in\{0,1\}^t}\zeta_{2^K}^{Q(x,y)}$ to be evaluated exactly as the algebraic product $\prod_{j=1}^{t}(1+\zeta_{2^K}^{c_j(y)})$, eliminating explicit enumeration of the $2^t$ superposition terms. 

We reinterpret the MINSPM principle of exploiting a small number of unique nonzero entries as an algebraic compression mechanism. Rather than explicitly enumerating the exponentially large set of superposition contributions, the proposed representation identifies the distinct nonzero phase contributions and encodes them in a compact dictionary whose values belong to a finite/Galois-ring-derived algebraic alphabet. A minimal indexing scheme provides direct access to the unique contributions and their associated multiplicities, replacing redundant enumeration by operations on the compressed representation. In the present generalized $\mathbb{Z}_{2^K}$-valued formulation, the separable phase structure permits the compressed sum to be evaluated analytically as a product of local factors in $\mathcal{O}(t)$ time. Thus, the method combines MINSPM-style minimal dictionary encoding with finite-ring algebraic factorization, transforming explicit superposition expansion into a compact exact representation while preserving the underlying amplitude algebra.
\end{abstract}

\section{Introduction}
The Quantum Fourier Transform (QFT) is the central algorithmic subroutine in exponentially advantageous quantum algorithms, most notably Shor's algorithm for integer factorization and the computation of discrete logarithms. However, the standard QFT circuit requires $\mathcal{O}(n^2)$ controlled-phase rotation gates and all-to-all qubit connectivity, imposing severe coherence and topological demands on physical quantum hardware. 

To mitigate these hardware constraints, Coppersmith introduced the Approximate Quantum Fourier Transform (AQFT), which truncates phase rotations below a threshold angle. By restricting the maximum coupling distance to a finite range $K$, the AQFT reduces the gate complexity to $\mathcal{O}(n \log n)$ while maintaining a bounded algorithmic error probability. Despite this hardware simplification, the classical simulation of the AQFT remains a formidable computational bottleneck.

Evaluating the output amplitude of a specific quantum state typically requires summing over the exponentially large basis of the quantum register. For a register of active size $t$, calculating a single Fourier amplitude involves evaluating a superposition of $2^t$ terms. Generalized classical simulation techniques, such as Matrix Product States (MPS) and tensor network contractions, evaluate these circuits using numerical floating-point tensor contractions, which do not inherently exploit the strict algebraic symmetries of the AQFT phase function.

In this paper, we resolve this simulation bottleneck for bounded-range AQFT circuits by introducing an exact finite-ring algebraic compression framework. We demonstrate that for any arbitrary coupling distance $K$, the AQFT phase function is strictly bilinear and completely separable in the superposition variables. This exact separability allows the exponentially large amplitude sum, previously requiring $\mathcal{O}(2^t)$ operations, to be analytically factored into a linear $\mathcal{O}(t)$ algebraic product over the ring $\mathbb{Z}_{2^K}$.

\section{Methodology: Universal Separability and Finite-Ring Factorization}

\subsection{The Generalized Coppersmith Phase Function}
We consider an AQFT circuit applied to a quantum register where the classical binary variables representing the superposition state are denoted by $x \in \{0, 1\}^t$, and the target measurement state is denoted by the fixed binary string $y \in \{0, 1\}^t$. 

Under Coppersmith's approximation with a maximum coupling range $K$, we map the phases into integer units of the minimum rotation angle $2\pi/2^K$. The resulting scalar phase function $Q(x,y)$, evaluated in $\mathbb{Z}_{2^K}$, is:
\begin{equation}
Q(x,y) = \sum_{i} 2^{K-1} x_i y_i + \sum_{j<i, i-j \le K} 2^{K-(i-j)} x_j y_i
\end{equation}

\subsection{Theorem: Universal Bilinear Separability}
\begin{theorem}
For any arbitrary AQFT truncation depth $K$, the resulting phase function $Q(x,y)$ is completely separable with respect to the active superposition variables $x_j$, containing no cross-coupling terms $x_j x_k$.
\end{theorem}

\begin{proof}
In the quantum circuit model of the QFT, controlled-phase gates entangle the input superposition register exclusively with the output target register. The total phase accumulated for any computational path is a sum of bipartite interactions between input bits $x_j$ and output bits $y_i$. Because we are evaluating the amplitude for a specific, fixed measurement state $y$, the variables $y_i \in \{0,1\}$ act as scalar constants. Consequently, the inner sum over the target index $i$ collapses into a single constant coefficient $c_j(y)$ for each superposition bit $x_j$. Thus, the phase function is strictly linear in $x$:
\begin{equation}
Q(x,y) = \sum_{j=1}^t c_j(y) x_j
\end{equation}
\end{proof}

\subsection{Generalized Finite-Ring Dictionary Mapping}
To evaluate the exact coefficients, we map the phase structure into the finite cyclotomic representation generated by the complex root $\zeta_{2^K} = e^{2\pi i/2^K}$. The phase coefficients map exactly to the ring of integers modulo $2^K$. 

By factoring $x_j$ out of the generalized phase function $Q(x,y)$, the coefficient for any specific variable $x_j$ is given by the closed-form integer equation:
\begin{equation}
c_j(y) = 2^{K-1} y_j + \sum_{i=j+1}^{j+K} 2^{K-(i-j)} y_i \pmod{2^K}
\end{equation}

\subsection{The MINSPM Factorization Identity}
Because the phase function $Q(x,y)$ is strictly separable for any $K$, the exponential sum over the $2^t$-dimensional state-space factors analytically. The total probability amplitude $A(y)$ corresponding to the measurement state $y$ evaluates exactly as:
\begin{equation}
\sum_{x\in\{0,1\}^t} \zeta_{2^K}^{Q(x,y)} = \sum_{x_1,\dots,x_t} \zeta_{2^K}^{\sum_j c_j(y) x_j} = \sum_{x_1,\dots,x_t} \prod_{j=1}^t \left(\zeta_{2^K}^{c_j(y)}\right)^{x_j}
\end{equation}

Summing independently over each binary variable $x_j \in \{0, 1\}$ yields the exact algebraic factorization identity:
\begin{equation}
A(y) = \prod_{j=1}^t \left( 1 + \zeta_{2^K}^{c_j(y)} \right)
\end{equation}

\section{Asymptotic Complexity Reduction}
The classical computation of the explicitly enumerated AQFT superposition scales as $\mathcal{O}(2^t)$, representing an exponential wall for state-vector simulators. 

By applying the finite-ring dictionary encoding, the evaluation of the total amplitude bypasses the sum entirely. Constructing the coefficient $c_j(y)$ for a specific bit requires a local summation over a finite window of size $K$. Performing this for all $t$ active variables, followed by evaluating the product $\prod (1 + \zeta_{2^K}^{c_j(y)})$, reduces the computational complexity to $\mathcal{O}(t \cdot K)$. Because the Coppersmith AQFT treats $K$ as a fixed, finite constant, the exact amplitude evaluation runs in strictly linear $\mathcal{O}(t)$ time.

\section{Conclusion}
We have established a mathematically rigorous framework for the exact analytical evaluation of Approximate Quantum Fourier Transform amplitudes. By proving the universal bilinear separability of the Coppersmith phase function, we demonstrated that AQFT superpositions can be losslessly compressed into a finite-ring dictionary over $\mathbb{Z}_{2^K}$. This yields an exact algebraic product that evaluates the amplitude of any target measurement state in $\mathcal{O}(t)$ time. This exact MINSPM algebraic factorization provides a foundation for classical verification and analysis of massively scaled quantum circuits without requiring tensor approximation or explicit superposition expansion.

\newpage
\appendix
\section{Generalized Validation Suite and Empirical Results}
To empirically verify the theorem of universal separability, we constructed a generalized software validation suite in the computational algebra system GAP. The suite evaluates arbitrary AQFT coupling ranges $K$, directly comparing the exponentially scaled classical Coppersmith expansion against the $\mathcal{O}(t)$ finite-ring MINSPM algebraic product.

\subsection{GAP Reference Implementation}
\begin{lstlisting}
# ===================================================================
# Generalized Coppersmith vs. MINSPM AQFT Validation Suite
# Includes High-K (K=10) and Large-t (1,048,576 terms) Benchmarks
# ===================================================================

Generalized_Z_Dictionary := function(phase_idx, K)
    local P;
    P := 2^K;
    return E(P)^(phase_idx mod P);
end;

# 1. Generalized Classical Coppersmith Scalar Phase
Coppersmith_Scalar_Gen := function(x_bits, y_bits, n_qubits, K)
    local phase_sum, i, j, k_diff, dist_factor, P;
    P := 2^K;
    phase_sum := 0;
    
    for i in [1..n_qubits] do
        if x_bits[i] = 1 and y_bits[i] = 1 then
            phase_sum := phase_sum + QuoInt(P, 2);
        fi;
        for j in [1..(i-1)] do
            k_diff := i - j;
            if k_diff <= K then
                if x_bits[j] = 1 and y_bits[i] = 1 then
                    dist_factor := QuoInt(P, 2^k_diff);
                    phase_sum := phase_sum + dist_factor;
                fi;
            fi;
        od;
    od;
    return Generalized_Z_Dictionary(phase_sum, K);
end;

# Brute-force summation for validation baseline O(2^t)
Compute_Coppersmith_Superposition_Gen := function(n_qubits, t_gates, y_vec, K)
    local num_terms, amp_sum, idx, x_vec, i, start_t, end_t;
    num_terms := 2^t_gates;
    amp_sum := 0 * E(2^K);
    
    start_t := Runtime();
    for idx in [1..num_terms] do
        x_vec := List([1..n_qubits], k -> 0);
        for i in [1..t_gates] do
            x_vec[i] := QuoInt(idx - 1, 2^(i-1)) mod 2;
        od;
        amp_sum := amp_sum + Coppersmith_Scalar_Gen(x_vec, y_vec, n_qubits, K);
    od;
    end_t := Runtime();
    
    return [amp_sum, (end_t - start_t) / 1000.0];
end;

# 2. Generalized MINSPM Finite-Ring Analytical Factorization O(t)
Compute_MINSPM_Algebraic_Gen := function(n_qubits, t_gates, y_vec, K)
    local total_amp, j, i, k_diff, dist_factor, coeff_j, P, start_t, end_t;
    P := 2^K;
    
    start_t := Runtime();
    total_amp := 1 * E(P)^0;
    
    for j in [1..t_gates] do
        coeff_j := 0;
        if y_vec[j] = 1 then
            coeff_j := coeff_j + QuoInt(P, 2);
        fi;
        
        for i in [(j+1)..n_qubits] do
            k_diff := i - j;
            if k_diff <= K then
                if y_vec[i] = 1 then
                    dist_factor := QuoInt(P, 2^k_diff);
                    coeff_j := coeff_j + dist_factor;
                fi;
            fi;
        od;
        
        total_amp := total_amp * (1 + Generalized_Z_Dictionary(coeff_j, K));
    od;
    end_t := Runtime();
    
    return [total_amp, (end_t - start_t) / 1000.0];
end;

# High-K Explicit and Randomized Test Harness
Run_High_K_Validation_Suite := function()
    local n, t, K, y_vec_explicit, y_rand, cop_res, min_res;
    
    Print("\n=================================================================\n");
    Print(" MINSPM vs COPPERSMITH: HIGH-K TRUNCATION & PERFORMANCE BENCHMARK\n");
    Print("=================================================================\n");
    
    # ---------------------------------------------------------
    # TEST 1: Explicit 64-qubit vector (K=10, t=20)
    # ---------------------------------------------------------
    n := 64;
    t := 20; 
    K := 10;
    
    y_vec_explicit := [
      1,0,1,1,0,0,1,0,
      1,1,0,1,0,1,1,0,
      0,1,0,1,1,0,1,1,
      1,0,0,1,0,1,0,1,
      1,1,0,0,1,1,0,1,
      0,1,1,0,1,0,0,1,
      1,0,1,1,0,1,1,0,
      0,0,1,0,1,1,0,1
    ];
    
    Print("\n--> TEST 1: Explicit Test Vector (n = 64, t = 20, K = 10, Ring Z_1024)\n");
    Print("    Evaluating 2^20 (1,048,576) explicit superposition terms...\n");
    
    cop_res := Compute_Coppersmith_Superposition_Gen(n, t, y_vec_explicit, K);
    min_res := Compute_MINSPM_Algebraic_Gen(n, t, y_vec_explicit, K);
    
    if cop_res[1] = min_res[1] then
        Print("    [PASS] Exact Algebraic Match Confirmed.\n");
    else
        Print("    [FAIL] Mismatch detected!\n");
    fi;
    
    Print("    - Classical Coppersmith O(2^t) Runtime : ", cop_res[2], " seconds\n");
    Print("    - MINSPM Algebraic O(t) Runtime        : ", min_res[2], " seconds\n");
    
    # ---------------------------------------------------------
    # TEST 2: Randomized Vector (K=12, t=20)
    # ---------------------------------------------------------
    K := 12;
    y_rand := List([1..n], i -> Random([0, 1]));
    
    Print("\n--> TEST 2: Randomized Test Vector (n = 64, t = 20, K = 12, Ring Z_4096)\n");
    Print("    Evaluating 2^20 (1,048,576) explicit superposition terms...\n");
    
    cop_res := Compute_Coppersmith_Superposition_Gen(n, t, y_rand, K);
    min_res := Compute_MINSPM_Algebraic_Gen(n, t, y_rand, K);
    
    if cop_res[1] = min_res[1] then
        Print("    [PASS] Exact Algebraic Match Confirmed.\n");
    else
        Print("    [FAIL] Mismatch detected!\n");
    fi;
    
    Print("    - Classical Coppersmith O(2^t) Runtime : ", cop_res[2], " seconds\n");
    Print("    - MINSPM Algebraic O(t) Runtime        : ", min_res[2], " seconds\n");
    Print("\n=================================================================\n");
end;

Run_High_K_Validation_Suite();
\end{lstlisting}

\end{document}
